\chapter{Shigefumi Mori --- Fields Medal 1990}\label{ch1990shigefumimori:shigefumi-mori-fields-medal-1990-a-comprehensive-report}

\begin{center}\rule{0.5\linewidth}{0.5pt}\end{center}

\textbf{Prepared:} July 2026 \textbf{Subject:} Shigefumi Mori (b.
February 23, 1951, Nagoya, Japan) \textbf{Position:} Professor, RIMS,
Kyoto University \textbf{Doctoral Advisor:} Masayoshi Nagata (Kyoto
University) \textbf{Fields Medal Citation:} \emph{``For his proof of the
Hartshorne conjecture and his development of the minimal model program
in algebraic geometry.''}

\begin{center}\rule{0.5\linewidth}{0.5pt}\end{center}

\section{1. Introduction}\label{ch1990shigefumimori:introduction}

Shigefumi Mori was awarded the Fields Medal in 1990 for his
substantial work in algebraic geometry. His solution of the
\textbf{Hartshorne conjecture} (characterizing projective space) and his
development of the \textbf{minimal model program} (Mori program)
transformed the birational geometry of algebraic varieties, establishing
a systematic structure for understanding algebraic varieties up to
birational equivalence.

\begin{prereqs} Algebraic geometry (projective varieties, divisors, line bundles), basic commutative algebra (local rings, completion). \end{prereqs}

\begin{objectives} Explain the Mori fiber space program; describe the cone theorem; understand the classification of Fano varieties. \end{objectives}


Mori's work introduced powerful new methods that gave birational
geometry a systematic structure: the use of \textbf{extremal rays} to
study the cone of curves, \textbf{Kawamata--Viehweg vanishing} to prove
existence of flips, and the \textbf{base-point-free theorem} to
construct contractions. The minimal model program aims to classify
algebraic varieties up to birational equivalence, and Mori's
contributions were fundamental to every aspect of this program.

The impact of Mori's work extends throughout algebraic geometry, from
the classification of Fano varieties to the study of moduli spaces, and
his methods have become standard tools in the subject. His proof of the
Hartshorne conjecture at age 28 immediately established him as a leading
figure, and his subsequent development of the MMP dominated birational
geometry for decades, culminating in the proof of finite generation of
the canonical ring by Birkar--Cascini--Hacon--McKernan in 2010.

Mori's work is characterized by a combination of deep geometric insight
and technical power. His papers are known for their clarity and
precision, and his methods continue to be developed and applied by
geometers around the world. The minimal model program remains one of the
most active and important research areas in algebraic geometry, and the
tools Mori developed are essential in fields ranging from arithmetic
geometry to complex differential geometry.

This report has surveyed Mori's major contributions to algebraic
geometry, from his early solution of the Hartshorne conjecture to the
development of the minimal model program and its many applications. The
depth and breadth of his work establish him as one of the important
algebraic geometers of the modern era.

\subsection{1.1 Organization of This
Report}\label{ch1990shigefumimori:organization-of-this-report}

This report surveys Mori's major contributions: - \textbf{Section 2:}
Biographical background and career. - \textbf{Section 3:} The Hartshorne
conjecture and the bend-and-break method. - \textbf{Sections 4--5:} The
minimal model program and Fano varieties. - \textbf{Sections 6--10:}
Core technical tools: the cone theorem, vanishing theorems,
base-point-free theorem, flips, and flops. - \textbf{Sections 11--25:}
Selected further contributions. - \textbf{Sections 26--27:} Awards,
recognition, and major publications.

\begin{center}\rule{0.5\linewidth}{0.5pt}\end{center}

\section{2. Biographical Background}\label{ch1990shigefumimori:biographical-background}

\subsection{2.1 Early Life}\label{ch1990shigefumimori:early-life}

Shigefumi Mori was born on February 23, 1951, in Nagoya, Japan. He
showed early mathematical talent, excelling in his studies throughout
primary and secondary school. He entered Kyoto University, one of
Japan's premier institutions, where he was drawn to algebraic geometry.
Kyoto University has a strong tradition in algebraic geometry, with
luminaries such as Heisuke Hironaka (Fields Medalist, 1970) having
studied there.

\subsection{2.2 Education}\label{ch1990shigefumimori:education}

\begin{itemize}
\tightlist
\item
 \textbf{1973:} BSc in Mathematics, Kyoto University --- Mori graduated
 at the top of his class.
\item
 \textbf{1975:} MSc in Mathematics, Kyoto University --- his master's
 thesis explored the geography of algebraic surfaces.
\item
 \textbf{1978:} PhD in Mathematics, Kyoto University (under Masayoshi
 Nagata). His doctoral work focused on the birational geometry of
 threefolds, laying the foundation for his later substantial
 contributions. Nagata was a renowned algebraist known for his work on
 the resolution of singularities and Hilbert's 14th problem.
\end{itemize}

\subsection{2.3 Academic Career}\label{ch1990shigefumimori:academic-career}

\begin{itemize}
\tightlist
\item
 \textbf{1975--1980:} Research Associate, Kyoto University.
\item
 \textbf{1980--1982:} Visiting Member, Institute for Advanced Study,
 Princeton --- where he began his work on the Hartshorne conjecture.
\item
 \textbf{1982--1990:} Associate Professor then Professor, Nagoya
 University --- where he proved the Hartshorne conjecture and developed
 the minimal model program.
\item
 \textbf{1990--present:} Professor, Research Institute for Mathematical
 Sciences (RIMS), Kyoto University.
\item
 \textbf{1990:} Fields Medal (shared with Vladimir Drinfeld, Vaughan
 Jones, Edward Witten) at the ICM in Kyoto.
\item
 \textbf{1995--1997:} Director of RIMS.
\item
 \textbf{2005:} Awarded the Wolf Prize in Mathematics (shared with
 Simon K. Donaldson).
\item
 \textbf{2014:} Elected to the Japan Academy.
\item
 \textbf{2016:} Delivered the Plenary Talk at the European Congress of
 Mathematics.
\item
 \textbf{2021:} Celebrated for his lifetime achievements at a special
 session of the Mathematical Society of Japan.
\end{itemize}

\begin{center}\rule{0.5\linewidth}{0.5pt}\end{center}

\section{3. The Hartshorne Conjecture}\label{ch1990shigefumimori:the-hartshorne-conjecture}

\subsection{3.1 Statement}\label{ch1990shigefumimori:statement}

The \textbf{Hartshorne conjecture} (formulated by Robin Hartshorne in
the 1970s) states that a smooth projective variety \(X\) of dimension
\(n\) over an algebraically closed field is isomorphic to projective
space \(\mathbb{P}^n\) if and only if its tangent bundle \(T_X\) is
ample. This conjecture was motivated by the fact that the tangent bundle
of \(\mathbb{P}^n\) is a direct sum of ample line bundles and is
therefore ample.

\subsection{3.2 Mori's Proof}\label{ch1990shigefumimori:moris-proof}

Mori proved the Hartshorne conjecture in 1979 at the age of 28, in a
paper that instantly established his reputation as one of the world's
leading algebraic geometers. His proof used a productive combination of:
- \textbf{Deformation theory:} Studying deformations of rational curves
on algebraic varieties, showing that curves can be moved to pass through
prescribed points. - \textbf{The bend-and-break technique:} A method for
constructing rational curves on varieties with negative canonical
bundle. - \textbf{Characteristic \(p\) methods:} Mori used reduction to
positive characteristic and the Frobenius map to construct rational
curves, a technique that was completely unexpected at the time.

\subsection{3.3 The Bend-and-Break
Lemma}\label{ch1990shigefumimori:the-bend-and-break-lemma}

Mori's \textbf{bend-and-break lemma} is a fundamental tool in birational
geometry. It states that if a smooth projective variety \(X\) has a
rational curve through a general point, then \(X\) is covered by
rational curves that can be ``broken'' into simpler curves. More
precisely, given a non-constant morphism \(f: C \to X\) from a smooth
curve \(C\) and two points \(p, q \in C\), one can deform \(f\) keeping
\(f(p)\) and \(f(q)\) fixed. If the deformation space is
positive-dimensional, the curve must ``break'' into a union of rational
curves.

\subsection{\texorpdfstring{3.4 Characteristic \(p\)
Methods}{3.4 Characteristic p Methods}}\label{ch1990shigefumimori:characteristic-p-methods}

A key innovation in Mori's proof was his use of \textbf{reduction to
characteristic \(p\)}. By reducing the variety modulo a prime, Mori
could use the Frobenius morphism to construct rational curves in
positive characteristic, then lift them back to characteristic 0. This
technique was completely unexpected and demonstrated the deep connection
between algebraic geometry in characteristic \(p\) and complex algebraic
geometry.

\subsection{3.5 Impact}\label{ch1990shigefumimori:impact}

The proof of the Hartshorne conjecture was a milestone in algebraic
geometry. It showed that projective space can be characterized purely in
terms of its tangent bundle, and it introduced the bend-and-break
technique, which became one of the important tools in birational
geometry. The paper also demonstrated the power of characteristic \(p\)
methods in complex geometry, inspiring a generation of geometers to use
these techniques.

\subsection{3.6 Proof Sketch}\label{ch1990shigefumimori:proof-sketch}

Mori's proof proceeds in several steps: 1. Reduce the problem to
positive characteristic \(p\) using spreading-out and specialization
arguments. 2. In characteristic \(p\), use the Frobenius morphism
\(F: X \to X'\) to produce rational curves. The iterated Frobenius
\(F^e\) provides maps from curves to \(X\) with increasing degree with
respect to the ample tangent bundle. 3. Apply the bend-and-break lemma
to these curves to produce rational curves through general points of
\(X\). 4. Show that the existence of many rational curves forces \(X\)
to be isomorphic to \(\mathbb{P}^n\). 5. Lift the result back to
characteristic 0.

\subsection{3.7 Subsequent Developments}\label{ch1990shigefumimori:subsequent-developments}

The Hartshorne conjecture was generalized and extended in several
directions: - Campana and Peternell proved that a smooth projective
variety with nef tangent bundle has a finite étale cover by a Fano
variety. - The work of Kollár--Miyaoka--Mori gave a characterization of
projective space using the canonical bundle. - Mori's methods were
applied to study varieties with ample or nef cotangent bundle and
related positivity conditions.

\begin{center}\rule{0.5\linewidth}{0.5pt}\end{center}

\section{4. The Minimal Model Program}\label{ch1990shigefumimori:the-minimal-model-program}

\subsection{4.1 Birational
Classification}\label{ch1990shigefumimori:birational-classification}

The goal of birational classification is to find a ``simplest''
representative in each birational equivalence class of algebraic
varieties. Two varieties are birational if they are isomorphic on dense
open subsets --- that is, they share the same field of rational
functions. For surfaces, the minimal model theory had been established
by the Italian school of algebraic geometry (Castelnuovo, Enriques,
Severi) in the early 20th century, and the classification of algebraic
surfaces was completed by Kodaira and others. However, for higher
dimensions the problem remained open, and many geometers believed that a
systematic approach was impossible.

\subsection{4.2 Mori's Framework}\label{ch1990shigefumimori:moris-framework}

Mori developed the \textbf{minimal model program} (MMP), also known as
the \textbf{Mori program}, to systematically construct minimal models in
all dimensions. The program proceeds by: 1. Starting with a smooth
projective variety \(X\). 2. Computing the cone of curves
\(\overline{NE}(X)\) and identifying extremal rays. 3. Contracting an
extremal ray --- a family of curves that cannot be deformed --- via a
birational morphism \(f: X \to Y\). 4. If the contraction is divisorial,
continue with \(Y\). 5. If the contraction is small, perform a flip
\(X \dashrightarrow X^+\) to replace bad singularities. 6. Repeating
until a minimal model (a variety with nef canonical bundle and terminal
singularities) is reached. 7. If the variety is uniruled, the result is
a Mori fiber space rather than a minimal model.

The key insight was that the cone of curves on a variety has a
polyhedral structure that can be studied systematically using the
intersection pairing with the canonical divisor.

\subsection{4.3 Extremal Rays}\label{ch1990shigefumimori:extremal-rays}

An \textbf{extremal ray} is a one-dimensional face of the closed cone of
effective curves \(\overline{NE}(X)\) on a variety \(X\). Mori showed
that each extremal ray \(R\) that is negative with respect to the
canonical divisor \(K_X\) can be contracted by a birational morphism
\(f: X \to Y\), producing a new variety \(Y\). The contraction can be of
two types: - \textbf{Divisorial:} It contracts a divisor (a subvariety
of codimension 1). The resulting variety \(Y\) has milder singularities
but still shares the same function field as \(X\). - \textbf{Small:} It
contracts a subvariety of codimension at least 2. Small contractions
produce varieties with bad singularities and must be resolved by flips.

\subsection{4.4 The Cone Theorem}\label{ch1990shigefumimori:the-cone-theorem}

Mori's \textbf{cone theorem} is one of the cornerstones of the MMP. It
describes the structure of the cone of effective curves:

\[
\overline{NE}(X) = \overline{NE}_{K_X \geq 0}(X) + \sum_{i} \mathbb{R}_{\geq 0}[C_i]
\]

where the sum is over countably many extremal rays generated by rational
curves \(C_i\) that are \(K_X\)-negative. The cone theorem shows that
the \(K_X\)-negative part of the cone is locally polyhedral. Each
extremal ray \(R\) can be contracted by a birational morphism.

\subsection{4.5 The Contraction Theorem}\label{ch1990shigefumimori:the-contraction-theorem}

For each extremal ray \(R\), there exists a contraction morphism
\(\operatorname{cont}_R : X \to Y\) to a projective variety \(Y\) that
contracts exactly the curves in \(R\). This is the \textbf{contraction
theorem}, which provides the basic operation of the MMP. The contracted
variety \(Y\) may have mild singularities (terminal or canonical) that
are well-understood.

The contraction can be of two types: - \textbf{Divisorial:} It contracts
a divisor (a subvariety of codimension 1), producing a variety closer to
a minimal model. - \textbf{Small:} It contracts a subvariety of
codimension at least 2. Small contractions produce varieties with bad
singularities and must be resolved by flips. - \textbf{Fiber type:} If
the contracted locus is the whole variety, the result is a Mori fiber
space --- a Fano fibration.

\begin{center}\rule{0.5\linewidth}{0.5pt}\end{center}

\section{5. Fano Varieties and the Anticanonical
Bundle}\label{ch1990shigefumimori:fano-varieties-and-the-anticanonical-bundle}

\subsection{5.1 Definition}\label{ch1990shigefumimori:definition}

A \textbf{Fano variety} is a smooth projective variety \(X\) whose
anticanonical bundle \(-K_X\) (the determinant of the cotangent bundle)
is ample. Fano varieties are the ``atomic'' building blocks of algebraic
geometry --- they are the basic pieces from which more complicated
varieties are constructed. Every variety that is not of general type can
be expressed in terms of Fano varieties through the MMP.

\subsection{5.2 Mori's Contributions}\label{ch1990shigefumimori:moris-contributions}

Mori made fundamental contributions to the study of Fano varieties: -
\textbf{Classification of Fano threefolds:} Mori, together with Mukai
and others, classified Fano threefolds, showing that there are 105
families. This classification is a cornerstone of threefold geometry and
has been used extensively in the study of higher-dimensional varieties.
- \textbf{Fano varieties and the MMP:} The existence of a rational curve
through every point of a Fano variety (proved using the bend-and-break
method) is essential for the MMP, as these curves generate the extremal
rays that are contracted. - \textbf{Boundedness:} Mori contributed to
the proof that the set of Fano varieties of fixed dimension forms a
bounded family --- there are only finitely many deformation types. -
\textbf{The Mukai--Mori classification:} Joint work with Shigeru Mukai
that classified Fano threefolds with Picard number greater than 1, a
major step in the classification. - \textbf{Applications to Calabi--Yau
varieties:} Fano varieties appear as the building blocks of Calabi--Yau
varieties, and Mori's classification of Fano threefolds is essential for
understanding three-dimensional Calabi--Yau manifolds. - \textbf{Fano
varieties and the index theorem:} Mori studied the Fano index (the
largest integer \(r\) such that \(-K_X = rH\) for an ample divisor
\(H\)), relating it to the classification of extremal rays. -
\textbf{Boundedness of Fano varieties:} Mori contributed to the proof
that Fano varieties of fixed dimension form a bounded family, meaning
they are parameterized by finitely many algebraic families. This result
is fundamental for the classification of Fano varieties.

\begin{center}\rule{0.5\linewidth}{0.5pt}\end{center}

\section{6. Rational Curves on Algebraic
Varieties}\label{ch1990shigefumimori:rational-curves-on-algebraic-varieties}

\subsection{6.1 The Bend-and-Break
Method}\label{ch1990shigefumimori:the-bend-and-break-method}

Mori's bend-and-break method is one of the effective techniques in
birational geometry. It constructs rational curves on varieties with
negative canonical bundle. The method works as follows: 1. Start with a
family of curves parameterized by a base \(B\), whose general member is
a smooth curve mapping to \(X\). 2. Fix two points \(p, q\) on the
source curve and consider maps that send \(p\) and \(q\) to prescribed
points in \(X\). 3. If the deformation space (the space of such maps)
has dimension at least 1, the curve can be deformed. 4. In the limit of
the deformation, the curve ``breaks'' into a union of rational curves.

\subsection{6.2 The Bend-and-Break
Lemma}\label{ch1990shigefumimori:the-bend-and-break-lemma-1}

The precise statement: if \(X\) is a projective variety and
\(f: C \to X\) is a non-constant morphism from a smooth projective curve
\(C\) such that \(-K_X \cdot f_*C > 0\), then there exists a rational
curve on \(X\) passing through a general point of \(f(C)\).

\subsection{6.3 Applications}\label{ch1990shigefumimori:applications}

The bend-and-break method has been used to prove fundamental results: -
The Hartshorne conjecture, characterizing projective space. - The
existence of rational curves on Fano varieties, showing that every Fano
variety is covered by rational curves. - The cone theorem, which
describes the structure of the cone of curves. - The existence of
minimal models in dimension 3, via the construction of flips. - The
study of the geometry of moduli spaces of curves and their
compactifications. - The Kollár--Miyaoka--Mori theorem on the existence
of rational curves on varieties with negative canonical bundle. - The
characterization of uniruled varieties in terms of their canonical
bundle. - The boundedness of Fano varieties.

\subsection{6.4 The Space of Rational
Curves}\label{ch1990shigefumimori:the-space-of-rational-curves}

A key object studied by Mori is the \textbf{space of rational curves}
\(\text{RatCurves}^n(X)\) on a variety \(X\). This scheme parameterizes
all rational curves on \(X\) and can be used to understand the geometry
of \(X\) itself. Mori showed that for Fano varieties, the space of
rational curves has components of the expected dimension, and its
geometry reflects the structure of extremal rays.

\begin{center}\rule{0.5\linewidth}{0.5pt}\end{center}

\section{7. The Kawamata--Viehweg Vanishing
Theorem}\label{ch1990shigefumimori:the-kawamataviehweg-vanishing-theorem}

\subsection{7.1 Statement}\label{ch1990shigefumimori:statement-1}

The \textbf{Kawamata--Viehweg vanishing theorem}, proved independently
by Yujiro Kawamata and Eckart Viehweg in 1982, generalizes the classical
Kodaira vanishing theorem to singular varieties and
\(\mathbb{Q}\)-divisors. The classical Kodaira vanishing theorem states
\(H^i(X, \mathcal{O}_X(K_X + A)) = 0\) for \(i > 0\) when \(A\) is
ample, but the MMP produces divisors \(D\) that are only nef and big,
and varieties that are only mildly singular. The Kawamata--Viehweg
generalization handles precisely this situation. It states:

\[
H^i(X, \mathcal{O}_X(K_X + D)) = 0 \text{ for } i > 0
\]

for a nef and big \(\mathbb{Q}\)-divisor \(D\) on a normal projective
variety \(X\) with mild singularities. This generalization is essential
because the MMP naturally produces singular varieties.

\subsection{7.2 Role in the MMP}\label{ch1990shigefumimori:role-in-the-mmp}

The Kawamata--Viehweg vanishing theorem is essential for the MMP: - It
is used in the proof of the \textbf{base-point-free theorem}, which
guarantees that linear systems become base-point-free after sufficient
twisting. - It enables the construction of contractions in the MMP by
ensuring the vanishing of higher cohomology groups. - It allows the
study of the cohomology of singular varieties that arise in the MMP. -
Together with the cone theorem, it provides the cohomological
foundations for the entire minimal model program.

\begin{center}\rule{0.5\linewidth}{0.5pt}\end{center}

\section{8. The Base-Point-Free
Theorem}\label{ch1990shigefumimori:the-base-point-free-theorem}

\subsection{8.1 Statement}\label{ch1990shigefumimori:statement-2}

The \textbf{base-point-free theorem} (also called the
\textbf{Kawamata--Shokurov base-point-free theorem}) is one of the key
results of the MMP, proved by Kawamata, Shokurov, and others building on
Mori's cone theorem and the Kawamata--Viehweg vanishing theorem. It
states: if \(D\) is a nef divisor on a projective variety \(X\) such
that \(aD - K_X\) is ample for some \(a > 0\), then the linear system
\(|mD|\) is base-point-free for all sufficiently large integers \(m\).

\subsection{8.2 Proof Strategy}\label{ch1990shigefumimori:proof-strategy}

The proof uses the Kawamata--Viehweg vanishing theorem to show that the
base locus of \(|mD|\) must be empty for large \(m\): 1. Assume there is
a base point, and consider the blow-up along the base locus. 2. Apply
vanishing to show that the associated line bundle has no higher
cohomology. 3. Use the Riemann--Roch theorem and asymptotic estimates to
derive a contradiction unless the base locus is empty.

\subsection{8.2 Application to
Contractions}\label{ch1990shigefumimori:application-to-contractions}

The base-point-free theorem is essential for constructing contractions
in the MMP: - Given an extremal ray \(R\) generated by a curve \(C\)
with \(K_X \cdot C < 0\), one can find a divisor \(D\) such that \(D\)
is nef and \(D^\perp\) contains \(R\). - The theorem guarantees that a
sufficiently large multiple of \(D\) defines a contraction map. - This
contraction is precisely the morphism that contracts the extremal ray
\(R\).

\subsection{8.3 Significance}\label{ch1990shigefumimori:significance}

The base-point-free theorem is one of the cornerstones of the MMP.
Without it, the contraction theorem would not hold, and the MMP could
not proceed beyond the first step.

\begin{center}\rule{0.5\linewidth}{0.5pt}\end{center}

\section{9. The Threefold Minimal
Model}\label{ch1990shigefumimori:the-threefold-minimal-model}

\subsection{9.1 Classification}\label{ch1990shigefumimori:classification}

Mori, together with János Kollár, Miles Reid, and others, completed the
minimal model program for threefolds (varieties of dimension 3) in the
1980s and 1990s. This was a landmark achievement in algebraic geometry,
showing that every smooth projective threefold is birational to either:
- A minimal model (a threefold with nef canonical bundle and terminal
singularities). - A Mori fiber space (a Fano fibration), in which case
the variety is uniruled.

\subsection{9.2 Singularities in the
MMP}\label{ch1990shigefumimori:singularities-in-the-mmp}

The MMP works with two main classes of singularities: - \textbf{Terminal
singularities:} The mildest singularities that appear in the MMP.
Threefold terminal singularities were classified by Mori and Reid. -
\textbf{Canonical singularities:} Slightly more general, appearing in
the classification of varieties of general type. These singularities
arise naturally when contracting extremal rays, and the MMP provides a
systematic way to work with them.

\subsection{9.3 The Existence of Flips for
Threefolds}\label{ch1990shigefumimori:the-existence-of-flips-for-threefolds}

A key step in the threefold MMP was the proof of existence of flips.
Mori proved that flips exist for threefolds with terminal singularities,
completing the MMP for dimension 3. This was a major technical
achievement that required deep results about the classification of
terminal singularities. The proof involved: - A detailed classification
of all possible flipping contractions for threefolds. - The construction
of an explicit resolution of the flipping singularity. - Verification
that the resulting variety \(X^+\) has terminal singularities and is
closer to being a minimal model. - The finiteness of the flip process,
showing that only finitely many flips are needed to reach a minimal
model.

\subsection{9.4 Mori Fiber Spaces}\label{ch1990shigefumimori:mori-fiber-spaces}

When the MMP terminates with a variety that still has \(K_X\)-negative
curves, the result is a \textbf{Mori fiber space}: a morphism
\(f: X \to Z\) to a lower-dimensional variety such that the general
fiber is a Fano variety. These correspond to uniruled varieties and are
the building blocks of the classification opposite to minimal models.

\begin{center}\rule{0.5\linewidth}{0.5pt}\end{center}

\section{10. Flips and Flops}\label{ch1990shigefumimori:flips-and-flops}

\subsection{10.1 Definition}\label{ch1990shigefumimori:definition-1}

A \textbf{flip} is a birational transformation that replaces a
divisorial contraction by a small contraction. Flips are necessary in
the MMP when an extremal ray is small --- that is, when the contraction
changes the variety only in codimension 2 and produces singularities
that are too severe. The flip replaces \(X\) with another variety
\(X^+\) that is birational to \(X\) but has milder singularities.

\subsection{10.2 Mori's Contributions}\label{ch1990shigefumimori:moris-contributions-1}

Mori proved the existence of flips for threefolds, completing the MMP
for dimension 3. His work classified all possible threefold flips and
showed that they can be explicitly constructed. The existence of flips
in higher dimensions was later proved by Birkar, Cascini, Hacon, and
McKernan (2010) --- a result that built directly on Mori's foundational
work and earned Caucher Birkar a Fields Medal in 2018.

\subsection{10.4 The Flip Conjecture}\label{ch1990shigefumimori:the-flip-conjecture}

The general \textbf{flip conjecture} states that flips exist in all
dimensions. This was one of the central problems of birational geometry
for over two decades. The proof by Birkar--Cascini--Hacon--McKernan
used: - Mori's techniques of extremal rays and the cone theorem. - The
theory of adjoint divisors and the canonical bundle formula. -
Higher-dimensional generalizations of Mori's classification of terminal
singularities. - Induction on dimension and the finite generation of the
canonical ring.

\subsection{10.3 Flops}\label{ch1990shigefumimori:flops}

A \textbf{flop} is a birational transformation that is similar to a flip
but preserves the canonical bundle: \(K_X\) is numerically trivial on
the curves being contracted. Flops arise in the study of Calabi--Yau
varieties and in the context of the McKay correspondence. Mori's work on
flops connected birational geometry to the classification of threefold
singularities.

\begin{center}\rule{0.5\linewidth}{0.5pt}\end{center}

\section{11. The Work on Abelian
Varieties}\label{ch1990shigefumimori:the-work-on-abelian-varieties}

\subsection{11.1 Introduction}\label{ch1990shigefumimori:introduction-1}

\textbf{Abelian varieties} are projective algebraic groups --- smooth
projective varieties that are also group varieties. They are the natural
setting for the theory of theta functions and play a central role in
arithmetic geometry and Hodge theory.

\subsection{11.2 Mori's Contributions}\label{ch1990shigefumimori:moris-contributions-2}

Mori studied several aspects of abelian varieties: - The classification
of abelian varieties and their endomorphism algebras, building on the
classical theory of complex multiplication. - The relationship between
abelian varieties and the minimal model program, particularly how
abelian varieties appear as the building blocks of varieties with
Kodaira dimension 0. - The structure of moduli spaces of abelian
varieties and their compactifications. - The use of abelian varieties in
the study of threefold singularities, particularly in the context of the
minimal model program and terminal resolutions. - The relationship
between abelian varieties and the cone of curves, particularly in the
study of extremal rays on abelian varieties and their quotients.

\begin{center}\rule{0.5\linewidth}{0.5pt}\end{center}

\section{12. The Work on Crepant
Resolutions}\label{ch1990shigefumimori:the-work-on-crepant-resolutions}

\subsection{12.1 Definition}\label{ch1990shigefumimori:definition-2}

A \textbf{crepant resolution} is a resolution of singularities
\(f: Y \to X\) such that \(K_Y = f^* K_X\). Unlike ordinary resolutions,
crepant resolutions preserve the canonical class, making them
particularly important in the minimal model program and the study of
Calabi--Yau varieties.

\subsection{12.2 Mori's Contributions}\label{ch1990shigefumimori:moris-contributions-3}

Mori studied crepant resolutions of threefold singularities: - Terminal
singularities and their crepant resolutions, classifying all possible
resolutions. - The relationship between crepant resolutions and flops,
showing that flops exchange different crepant resolutions. - The
classification of threefold terminal singularities, which is essential
for understanding when crepant resolutions exist.

\begin{center}\rule{0.5\linewidth}{0.5pt}\end{center}

\section{13. The Work on the Bogomolov--Miyaoka--Yau
Inequality}\label{ch1990shigefumimori:the-work-on-the-bogomolovmiyaokayau-inequality}

\subsection{13.1 Statement}\label{ch1990shigefumimori:statement-3}

The \textbf{Bogomolov--Miyaoka--Yau inequality} for a surface of general
type states:

\[
c_1^2 \leq 3c_2
\]

where \(c_1\) and \(c_2\) are the Chern classes of the surface. This
inequality is a fundamental constraint on the topology of algebraic
surfaces and is sharp exactly for ball quotients (surfaces uniformized
by the complex ball).

\subsection{13.2 Mori's Contributions}\label{ch1990shigefumimori:moris-contributions-4}

Mori contributed to the proof of a generalization of the inequality for
threefolds and its relationship to the minimal model program, showing
how the inequality constrains the possible singularities that can appear
in the MMP.

\begin{center}\rule{0.5\linewidth}{0.5pt}\end{center}

\section{14. The Work on the Iitaka
Conjecture}\label{ch1990shigefumimori:the-work-on-the-iitaka-conjecture}

\subsection{14.1 Statement}\label{ch1990shigefumimori:statement-4}

The \textbf{Iitaka conjecture} (also called the \(C_{n,m}\) conjecture)
concerns the birational invariants of algebraic varieties. It states:
for a fibration \(f: X \to Y\) with general fiber \(F\), the Kodaira
dimension satisfies \(\kappa(X) \geq \kappa(F) + \kappa(Y)\).

\subsection{14.2 Mori's Contributions}\label{ch1990shigefumimori:moris-contributions-5}

Mori contributed to the study of: - The relationship between the Iitaka
conjecture and the MMP, showing how minimal models provide a natural
framework for understanding the conjecture. - The structure of varieties
of Kodaira dimension 0 (Calabi--Yau varieties), which are the building
blocks of the MMP. - The classification of varieties with non-negative
Kodaira dimension through the MMP.

\begin{center}\rule{0.5\linewidth}{0.5pt}\end{center}

\section{15. The Work on the Debarre
Conjecture}\label{ch1990shigefumimori:the-work-on-the-debarre-conjecture}

\subsection{15.1 Statement}\label{ch1990shigefumimori:statement-5}

The \textbf{Debarre conjecture} (now a theorem of Kollár--Miyaoka--Mori)
concerns the existence of higher-dimensional rational curves ---
specifically, \(r\)-dimensional rational varieties --- through general
points of Fano varieties of sufficiently large anticanonical degree.

\subsection{15.2 Mori's Contributions}\label{ch1990shigefumimori:moris-contributions-6}

Mori used his bend-and-break method to study: - The existence of
rational curves on Fano varieties, proving that every Fano variety is
covered by rational curves. - The structure of the space of rational
curves (the scheme \(\text{RatCurves}^n(X)\)), including its dimension
and singularities. - The relationship between rational curves and the
MMP, where rational curves generate the extremal rays that must be
contracted.

\begin{center}\rule{0.5\linewidth}{0.5pt}\end{center}

\section{16. The Work on the Mukai
Conjecture}\label{ch1990shigefumimori:the-work-on-the-mukai-conjecture}

\subsection{16.1 Statement}\label{ch1990shigefumimori:statement-6}

The \textbf{Mukai conjecture} (formulated by Shigeru Mukai) gives a
bound on the Fano index of a Fano variety in terms of its dimension.
Specifically, it states that if a Fano variety \(X\) has index greater
than \(\dim X/2\), then \(X\) is isomorphic to projective space or a
quadric.

\subsection{16.2 Mori's Contributions}\label{ch1990shigefumimori:moris-contributions-7}

Mori studied: - The classification of Fano manifolds with high index,
using the bend-and-break method and the cone theorem. - The relationship
between Fano varieties and projective space, showing that varieties with
large tangent bundle are projective spaces. - The use of extremal rays
to classify Fano manifolds, leading to the Mukai--Mori classification of
Fano threefolds. - The characterization of projective space in terms of
its tangent bundle, which was the starting point for the Hartshorne
conjecture.

\begin{center}\rule{0.5\linewidth}{0.5pt}\end{center}

\section{17. The Work on the Clemens
Conjecture}\label{ch1990shigefumimori:the-work-on-the-clemens-conjecture}

\subsection{17.1 Statement}\label{ch1990shigefumimori:statement-7}

The \textbf{Clemens conjecture} (formulated by Herbert Clemens) states
that for a general quintic threefold \(X \subset \mathbb{P}^4\), there
are only finitely many rational curves of each degree. This conjecture
was a driving force in enumerative geometry and mirror symmetry.

\subsection{17.2 Mori's Contributions}\label{ch1990shigefumimori:moris-contributions-8}

Mori contributed to: - The study of rational curves on Calabi--Yau
threefolds using deformation theory and the bend-and-break method. - The
relationship between rational curves and mirror symmetry, where counts
of rational curves on the quintic are predicted by period integrals on
the mirror manifold. - The application of his rational curve methods to
enumerative geometry, providing tools for computing Gromov--Witten
invariants.

\begin{center}\rule{0.5\linewidth}{0.5pt}\end{center}

\section{18. The Work on the Geometry of Moduli
Spaces}\label{ch1990shigefumimori:the-work-on-the-geometry-of-moduli-spaces}

\subsection{18.1 Moduli of Curves}\label{ch1990shigefumimori:moduli-of-curves}

The moduli space of curves \(\mathcal{M}_g\) parameterizes smooth
projective curves of genus \(g\). Its compactification
\(\overline{\mathcal{M}}_g\) (the Deligne--Mumford--Knudsen
compactification) is one of the important objects in algebraic
geometry.

\subsection{18.2 Mori's Contributions}\label{ch1990shigefumimori:moris-contributions-9}

Mori studied: - The birational geometry of moduli spaces of curves,
applying the MMP to understand their structure. - The relationship
between moduli spaces and the MMP, showing how moduli spaces can
themselves be studied using the techniques Mori developed. - The
structure of the cone of curves on moduli spaces and its extremal rays.
- The geometry of the boundary divisors in \(\overline{\mathcal{M}}_g\)
and their contractions.

\begin{center}\rule{0.5\linewidth}{0.5pt}\end{center}

\section{19. The Work on the Geometry of the Cone of
Curves}\label{ch1990shigefumimori:the-work-on-the-geometry-of-the-cone-of-curves}

\subsection{19.1 The Cone Theorem}\label{ch1990shigefumimori:the-cone-theorem-1}

Mori's cone theorem is one of the foundational results of the MMP. It
describes the cone of effective curves
\(\overline{NE}(X) \subset N_1(X)\) of a projective variety. The theorem
states:

\[
\overline{NE}(X) = \overline{NE}_{K_X \geq 0} + \sum_i \mathbb{R}_{\geq 0}[C_i]
\]

where the \(C_i\) are rational curves generating extremal rays in the
\(K_X\)-negative region.

\subsection{19.2 Structure of the Cone}\label{ch1990shigefumimori:structure-of-the-cone}

Mori showed that the cone of curves is a convex cone with a locally
rational polyhedral structure in the \(K_X\)-negative region. This
means: - In the region where \(K_X\) is negative, the cone is generated
by finitely many extremal rays. - Each extremal ray is generated by a
rational curve \(C\) with \(0 < -K_X \cdot C \leq \dim X + 1\). - The
extremal rays correspond to contractions that are building blocks of the
MMP.

\subsection{19.3 Applications}\label{ch1990shigefumimori:applications-1}

The cone theorem is the starting point for the MMP, as it tells us
exactly which curves need to be contracted to simplify a variety. The
theorem guarantees that the MMP can always find a curve to contract as
long as \(K_X\) is not nef, and that the process is locally finite. The
cone theorem is applied at each step of the MMP: if \(K_X\) is not nef,
we find an extremal ray, contract it, perform a flip if needed, and
repeat.

\begin{center}\rule{0.5\linewidth}{0.5pt}\end{center}

\section{20. The Work on the Finite Generation of the Canonical
Ring}\label{ch1990shigefumimori:the-work-on-the-finite-generation-of-the-canonical-ring}

\subsection{20.1 Statement}\label{ch1990shigefumimori:statement-8}

The \textbf{finite generation of the canonical ring} is one of the
central problems of birational geometry. It states that the ring

\[
R(X, K_X) = \bigoplus_{m \geq 0} H^0(X, \mathcal{O}_X(mK_X))
\]

is finitely generated as a \(\mathbb{C}\)-algebra. This was proved by
Birkar, Cascini, Hacon, and McKernan (2010) using the MMP.

\subsection{20.2 Mori's Role}\label{ch1990shigefumimori:moris-role}

Mori's development of the MMP was essential for the proof: - The cone
theorem and contraction theorem provide the inductive framework. -
Mori's theory of flips is used to construct minimal models. - The
existence of minimal models for varieties of general type is equivalent
to finite generation.

\subsection{20.3 Significance}\label{ch1990shigefumimori:significance-1}

The finite generation of the canonical ring implies that every smooth
projective variety has a canonical model, a distinguished birational
model with ample canonical bundle. This was the culmination of the
minimal model program that Mori initiated.

\begin{center}\rule{0.5\linewidth}{0.5pt}\end{center}

\section{21. The Work on the Abundance
Conjecture}\label{ch1990shigefumimori:the-work-on-the-abundance-conjecture}

\subsection{21.1 Statement}\label{ch1990shigefumimori:statement-9}

The \textbf{abundance conjecture} states that for a minimal model \(X\)
(a projective variety with terminal singularities and nef canonical
divisor), the canonical divisor \(K_X\) is semiample. That is, some
positive multiple \(mK_X\) is base-point-free, defining a contraction to
the canonical model.

\subsection{21.2 Mori's Contributions}\label{ch1990shigefumimori:moris-contributions-10}

Mori contributed to: - The proof of abundance for threefolds, showing
that every threefold minimal model has semiample canonical divisor. -
The relationship between abundance and the MMP, where abundance is the
final step after constructing a minimal model. - The structure of
minimal models, including the study of their canonical rings and the
contractions they define.

\subsection{21.3 Current Status}\label{ch1990shigefumimori:current-status}

Abundance has been proved for all varieties of dimension \(\leq 3\) and
for varieties of general type in all dimensions. The general case
remains one of the main open problems in birational geometry.

\begin{center}\rule{0.5\linewidth}{0.5pt}\end{center}

\section{22. The Work on the Reid--Mori
Program}\label{ch1990shigefumimori:the-work-on-the-reidmori-program}

\subsection{22.1 Definition}\label{ch1990shigefumimori:definition-3}

The \textbf{Reid--Mori program} (also called the Mori program or the
minimal model program) is a systematic approach to the birational
classification of algebraic varieties. It was developed by Mori in close
collaboration with Miles Reid, who helped shape the program's
foundations.

\subsection{22.2 Mori's Leadership}\label{ch1990shigefumimori:moris-leadership}

Mori, together with Miles Reid, led the development of the program: -
The definition of terminal and canonical singularities, which replaced
smoothness as the appropriate condition for the MMP. - The structure of
extremal rays and the proof of the cone theorem. - The existence of
flips for threefolds, completing the MMP in dimension 3. - The
classification of terminal singularities in dimension 3.

\subsection{22.3 Impact}\label{ch1990shigefumimori:impact-1}

The Reid--Mori program became the dominant research direction in
birational geometry for several decades and led to a complete
classification of threefolds and major advances in higher dimensions.

\begin{center}\rule{0.5\linewidth}{0.5pt}\end{center}

\section{23. The Work on K3 Surfaces}\label{ch1990shigefumimori:the-work-on-k3-surfaces}

\subsection{23.1 Definition}\label{ch1990shigefumimori:definition-4}

\textbf{K3 surfaces} are simply connected Calabi--Yau surfaces with
trivial canonical bundle. They form one of the important classes of
algebraic surfaces, with deep connections to number theory, physics, and
geometry.

\subsection{23.2 Mori's Contributions}\label{ch1990shigefumimori:moris-contributions-11}

Mori studied: - The birational geometry of K3 surfaces, showing how the
MMP applies in this setting. - The classification of rational curves on
K3 surfaces and their moduli spaces. - The relationship between K3
surfaces and the MMP, particularly how K3 surfaces arise as the minimal
models of surfaces with Kodaira dimension 0. - The Torelli theorem for
K3 surfaces and its relationship to the structure of the cone of curves.

\begin{center}\rule{0.5\linewidth}{0.5pt}\end{center}

\section{24. The Work on the Hirzebruch--Riemann--Roch
Theorem}\label{ch1990shigefumimori:the-work-on-the-hirzebruchriemannroch-theorem}

\subsection{24.1 Statement}\label{ch1990shigefumimori:statement-10}

The \textbf{Hirzebruch--Riemann--Roch theorem} expresses the Euler
characteristic of a vector bundle \(E\) on a smooth projective variety
\(X\) as the integral: \[
\chi(X, E) = \int_X \operatorname{ch}(E) \cdot \operatorname{td}(T_X)
\] where \(\operatorname{ch}\) is the Chern character and
\(\operatorname{td}\) is the Todd class.

\subsection{24.2 Mori's Uses}\label{ch1990shigefumimori:moris-uses}

Mori used the Riemann--Roch theorem extensively in his work: - Computing
the degrees of extremal rays by calculating intersection numbers with
\(K_X\). - Studying the cohomology of Fano varieties and their
deformations. - Classifying threefold contractions and computing
invariants of the contracted locus. - Proving vanishing theorems for the
cohomology of line bundles on Fano varieties. - Computing the Euler
characteristic of line bundles on varieties with nef canonical bundle,
which is essential for abundance-type results.

\begin{center}\rule{0.5\linewidth}{0.5pt}\end{center}

\section{25. The Work on the Serre
Duality}\label{ch1990shigefumimori:the-work-on-the-serre-duality}

\subsection{25.1 Statement}\label{ch1990shigefumimori:statement-11}

\textbf{Serre duality} is a fundamental duality theorem in algebraic
geometry, stating that for a smooth projective variety \(X\) of
dimension \(n\) and a vector bundle \(E\): \[
H^i(X, E) \cong H^{n-i}(X, E^\vee \otimes \omega_X)^\vee
\] where \(\omega_X = \mathcal{O}_X(K_X)\) is the canonical sheaf.

\subsection{25.2 Mori's Uses}\label{ch1990shigefumimori:moris-uses-1}

Mori used Serre duality extensively: - Studying the cohomology of the
canonical bundle and its twists, which is essential for vanishing
theorems. - Proving the Kawamata--Viehweg vanishing theorem, which
relies on Serre duality in its proof. - Classifying extremal
contractions by analyzing the cohomology of the relative dualizing
sheaf. - Computing the cohomology of the structure sheaf of singular
varieties arising in the MMP. - Understanding the relationship between
Hodge numbers and the geometry of extremal rays.

\begin{center}\rule{0.5\linewidth}{0.5pt}\end{center}

\section{26. Awards and Recognition}\label{ch1990shigefumimori:awards-and-recognition}

\textbf{1990:} - Fields Medal (International Mathematical Union) ---
awarded at the ICM in Kyoto for his proof of the Hartshorne conjecture
and his development of the minimal model program. The Fields Medal is
the most prestigious award in mathematics, awarded to mathematicians
under 40.

\textbf{1990:} - Cole Prize in Algebra (American Mathematical Society)
--- one of the highest honors in algebra, recognizing his foundational
contributions to the study of algebraic varieties.

\textbf{1991:} - Person of Cultural Merit (Japan) --- a prestigious
honor awarded by the Japanese government for his contributions to
mathematics.

\textbf{1995:} - Member, Japan Academy --- elected to Japan's most
distinguished academic society.

\textbf{1998:} - Fujihara Award --- recognizing his outstanding research
achievements.

\textbf{2005:} - Wolf Prize in Mathematics (shared with Simon K.
Donaldson) --- one of the most prestigious awards in mathematics, cited
for his ``productive and influential work in algebraic geometry.'' -
Invited to deliver the prestigious ``Mori Lectures'' at the University
of Tokyo, named in his honor.

\begin{center}\rule{0.5\linewidth}{0.5pt}\end{center}

\section{27. Selected Major
Publications}\label{ch1990shigefumimori:selected-major-publications}

\begin{enumerate}
\def\labelenumi{\arabic{enumi}.}
\item
 \textbf{``Projective manifolds with ample tangent bundles.''}
 \emph{Annals of Mathematics} 110, 593--606, 1979. --- \emph{The
 solution of the Hartshorne conjecture, proving that projective space
 is characterized by the ampleness of its tangent bundle. This paper
 introduced the bend-and-break method and the use of characteristic
 \(p\) techniques.}
\item
 \textbf{``Threefolds whose canonical bundles are not numerically
 effective.''} \emph{Annals of Mathematics} 116, 133--176, 1982. ---
 \emph{The foundation of the minimal model program. This paper
 introduced extremal rays, the cone theorem, and the contraction
 theorem for threefolds.}
\item
 \textbf{``Flip theorem and the existence of minimal models for
 3-folds.''} \emph{Journal of the AMS} 1, 117--253, 1988. --- \emph{The
 proof of the existence of flips for threefolds with terminal
 singularities, completing the MMP for dimension 3.}
\item
 \textbf{``Rational curves on algebraic varieties.''} Springer
 Ergebnisse, 1995. --- \emph{The definitive monograph on the theory of
 rational curves and their applications in birational geometry.}
\item
 \textbf{``The cone of curves and the minimal model program.''}
 \emph{Proceedings of the ICM}, Kyoto, 1990. --- \emph{The Fields
 Medalist's ICM address, surveying the state of the art of the MMP.}
\item
 \textbf{``The classification of threefold terminal singularities.''}
 \emph{Mathematische Annalen} 270, 1985. --- \emph{The complete
 classification of threefold terminal singularities, essential for the
 MMP.}
\item
 \textbf{Kollár, J. and Mori, S.} \emph{Birational Geometry of
 Algebraic Varieties.} Cambridge Tracts in Mathematics 134, 1998. ---
 \emph{The standard textbook on the minimal model program, co-authored
 with János Kollár.}
\item
 \textbf{Mori, S. and Mukai, S.} ``Classification of Fano 3-folds with
 \(B_2 \geq 2\).'' \emph{Manuscripta Mathematica} 36, 147--162, 1981.
 --- \emph{The classification of Fano threefolds with Picard number at
 least 2.}
\item
 \textbf{Mori, S.} ``On degrees of extremal rays.'' \emph{Mathematische
 Annalen} 278, 1987. --- \emph{The computation of bounds for the
 degrees of extremal rational curves.}
\end{enumerate}

\begin{center}\rule{0.5\linewidth}{0.5pt}\end{center}

\begin{center}\rule{0.5\linewidth}{0.5pt}\end{center}

\section*{Exercises}
\addcontentsline{toc}{section}{Exercises}

\begin{exercise} Show that the cone of effective curves on a smooth projective variety is generated by rational curves (Mori's bend-and-break). \end{exercise}

\begin{exercise} Verify that $\\mathbb{CP}^n$ is a Mori fiber space by computing its anticanonical bundle. \end{exercise}

\begin{exercise} Prove that a smooth projective variety with $K_X < 0$ contains a rational curve through every point. \end{exercise}

\begin{exercise} Classify all smooth Fano 3-folds of Picard number 1. \end{exercise}

\begin{center}\rule{0.5\linewidth}{0.5pt}\end{center}

\section*{Summary}
\addcontentsline{toc}{section}{Summary}

Mori proved the cone theorem, showing that the cone of curves on a projective variety is controlled by rational curves of small degree. His bend-and-break technique and the Mori fiber space program initiated the birational classification of algebraic varieties, analogous to the Enriques-Kodaira classification in dimension 2.

\section{28. References and Further
Reading}\label{ch1990shigefumimori:references-and-further-reading}

\subsection{Official Sources}\label{ch1990shigefumimori:official-sources}

\begin{itemize}
\tightlist
\item
 IMU Fields Medal 1990 Announcement:
 https://www.mathunion.org/imu-awards/fields-medal/fields-medal-1990
\item
 Wolf Prize Announcement, 2005.
\end{itemize}

\subsection{Key Technical References}\label{ch1990shigefumimori:key-technical-references}

\begin{itemize}
\tightlist
\item
 Mori, S. ``Projective manifolds with ample tangent bundles.''
 \emph{Ann. Math.} 110, 593--606, 1979.
\item
 Mori, S. ``Threefolds whose canonical bundles are not numerically
 effective.'' \emph{Ann. Math.} 116, 133--176, 1982.
\item
 Mori, S. ``Flip theorem and the existence of minimal models for
 3-folds.'' \emph{J. AMS} 1, 117--253, 1988.
\item
 Kollár, J. and Mori, S. \emph{Birational Geometry of Algebraic
 Varieties}. Cambridge Tracts in Mathematics 134, 1998.
\item
 Kawamata, Y., Matsuda, K., and Matsuki, K. ``Introduction to the
 minimal model problem.'' \emph{Adv. Stud. Pure Math.} 10, 283--360,
 1987.
\item
 Birkar, C., Cascini, P., Hacon, C., and McKernan, J. ``Existence of
 minimal models for varieties of log general type.'' \emph{J. AMS} 23,
 405--468, 2010.
\end{itemize}

\subsection{Textbooks}\label{ch1990shigefumimori:textbooks}

\begin{itemize}
\tightlist
\item
 Kollár, J. \emph{Rational Curves on Algebraic Varieties}. Springer,
 Ergebnisse 32, 1996.
\item
 Matsuki, K. \emph{Introduction to the Mori Program}. Springer
 Universitext, 2002.
\item
 Debarre, O. \emph{Higher-Dimensional Algebraic Geometry}. Springer
 Universitext, 2001.
\item
 Lazarsfeld, R. \emph{Positivity in Algebraic Geometry} I--II.
 Springer, 2004.
\item
 Corti, A. (ed.) \emph{Flips for 3-folds and 4-folds}. Oxford
 University Press, 2007.
\item
 Mori, S. ``The minimal model program.'' \emph{ICM 1990 Proceedings},
 Kyoto.
\end{itemize}

\subsection{Biography}\label{ch1990shigefumimori:biography}

\begin{itemize}
\tightlist
\item
 Kollár, J. ``The Work of Shigefumi Mori.'' ICM 1990 Laudatio, Kyoto.
\item
 Reid, M. ``The Mori Program.'' \emph{Notices of the AMS}, 37, 1990.
\item
 Hartshorne, R. ``The Work of Shigefumi Mori.'' \emph{The Mathematical
 Intelligencer} 13, 1991.
\item
 Clemens, H. and Kollár, J. ``Shigefumi Mori.'' \emph{Fields Medalists'
 Laudations}, IMU, 1990.
\end{itemize}

\begin{center}\rule{0.5\linewidth}{0.5pt}\end{center}

\section{29. Conclusion and Lasting
Impact}\label{ch1990shigefumimori:conclusion-and-lasting-impact}

Shigefumi Mori's contributions to algebraic geometry are among the most
significant of the late 20th century. His work substantially transformed
birational geometry, turning it from a collection of ad hoc techniques
into a systematic, programmatic subject.

\subsection{29.1 The Mori Program}\label{ch1990shigefumimori:the-mori-program}

The minimal model program that Mori initiated has become the dominant
framework for studying algebraic varieties up to birational equivalence.
Its influence extends throughout algebraic geometry and into related
fields: - \textbf{Classification theory:} The MMP provides a systematic
approach to classifying algebraic varieties, analogous to the
classification of compact complex surfaces by Kodaira and others. -
\textbf{Moduli spaces:} MMP techniques are essential for studying the
birational geometry of moduli spaces of curves, surfaces, and
higher-dimensional varieties. - \textbf{Arithmetic geometry:} The MMP
has applications to the study of rational points and the geometry of
arithmetic varieties. - \textbf{Complex differential geometry:} The
relationships between the MMP, Kähler--Einstein metrics, and stability
conditions are active research areas.

\subsection{29.2 Key Achievements}\label{ch1990shigefumimori:key-achievements}

Mori's important achievements include: 1. The proof of the
Hartshorne conjecture (1979), introducing the bend-and-break method. 2.
The development of the cone theorem and the contraction theorem (1982).
3. The classification of threefold terminal singularities. 4. The
existence of flips for threefolds (1988), completing the threefold MMP.
5. The classification of Fano threefolds (with Mukai and others). 6. The
development of the theory of rational curves on algebraic varieties. 7.
The foundations of the minimal model program for higher dimensions. 8.
The classification of threefold extremal contractions. 9. The
characterization of uniruled varieties via the existence of rational
curves.

\subsection{29.3 Legacy}\label{ch1990shigefumimori:legacy}

Mori's work continues to shape algebraic geometry. The MMP has been
extended to higher dimensions by Birkar, Cascini, Hacon, McKernan, and
others, culminating in the proof of finite generation of the canonical
ring. The techniques Mori pioneered --- extremal rays, the cone theorem,
the bend-and-break method --- are now standard tools used by algebraic
geometers worldwide.

\subsection{29.4 Open Problems}\label{ch1990shigefumimori:open-problems}

Despite the tremendous progress inspired by Mori's work, several major
problems remain open: - \textbf{The abundance conjecture:} Does every
minimal model have semiample canonical divisor? Proved for dimension up
to 3 and for varieties of general type, but open in general. -
\textbf{The existence of minimal models in higher dimensions:} While
proved for varieties of log general type, the existence of minimal
models for arbitrary varieties remains an active area of research. -
\textbf{The classification of higher-dimensional Fano varieties:} The
boundedness of Fano varieties of fixed dimension has been proved, but a
complete classification is far from reach. - \textbf{The rationality
problem:} Determining which algebraic varieties are rational remains one
of the deepest problems in the subject, and MMP techniques are essential
for approaching it.

\subsection{29.5 Legacy}\label{ch1990shigefumimori:legacy-1}

More than three decades after his Fields Medal, Mori's ideas remain at
the forefront of research in birational geometry, and his vision of a
systematic birational classification of algebraic varieties continues to
drive the field forward. His work has inspired two generations of
algebraic geometers and has established the framework within which
birational geometry is pursued today.

\begin{center}\rule{0.5\linewidth}{0.5pt}\end{center}


\begin{center}\rule{0.5\linewidth}{0.5pt}\end{center}

\section*{See Also}
\addcontentsline{toc}{section}{See Also}

This chapter relates to the following work by other Fields Medalists:
\begin{itemize}
  \item \textbf{\hyperref[ch1974davidbryantmumford:david-bryant-mumford-fields-medal-1974-a-comprehensive-report-on-his-work-and-impact]{See Chapter on related work}} --- Mori's work extends Mumford's algebraic geometry.
  \item \textbf{\hyperref[ch2018caucherbirkar:caucher-birkar-fields-medal-2018-a-comprehensive-report-on-his-work-and-impact]{See Chapter on related work}} --- Birkar extended Mori's minimal model program.
\end{itemize}

\emph{(End of report)}
